\chapter{先手算十一元五平均，再看一般奇数余量}\label{ch:odd}
\section{现在要做的，是把任务单上的资源造出来}
前面已经知道足够的剪切和同余子群能完成核心。现在不能继续把它们当作假设，必须给出真实位置选择。本讲取 $n=2p+r$，$r$ 为奇数且 $1\le r<p$，包括最紧的 $r=1$。

为方便交换两等大块，用对称坐标
\[u=z,qquad v=a-rz/2,qquad
A=ru/2+v,quad B=ru/2-v,quad C=-pu.
\]
交换 $A,B$ 恰为 $S=\diag(1,-1)$。这就是引入二分之一的目的：把一个原来稍复杂的角色交换变成直接给第二坐标变号。

\section{先看p等于五、r等于一}
从五份 $A$、五份 $B$、一份 $C$ 出发。
\begin{enumerate}
\item 先拿两份 $A$、三份 $B$ 平均，得到五份新值 $D$。
\item 保留一份 $D$ 作最终单点。剩下 $(A,B,C,D)$ 四类库存为 $(3,2,1,4)$。
\item 取输出组 $(2,1,0,2)$，即两份 $A$、一份 $B$、两份 $D$，做五平均。再把余下五个位置平均。
\end{enumerate}
做完又是两个五块加一个单点。这是一条三步返回。

另一条返回先拿四份 $A$、一份 $B$，保留一份新值后库存为 $(1,4,1,4)$，输出取 $(1,2,0,2)$，再平均补集。注意，两套表是\emph{可选的两条完整词}，不是要求同一轮同时取用两套库存。

在 $(u,v)$ 坐标里，两张清去共同分母的矩阵为
\[
J_{-1}=\begin{pmatrix}-5/2&1\\55/4&5/2\end{pmatrix},qquad
J_3=\begin{pmatrix}-5/2&-3\\55/4&5/2\end{pmatrix}.
\]
实际尺度为 $1/25$，两者平方为 $20I,-35I$。因此两条返回都取得了正向射影逆。两者的商为三角矩阵，倍率为 $-7/4$；与 $S$ 作交换子恰给 $U(1)$。倍率及其逆生成 $R=\Z[1/14]$，因而得到 $U(R)$。

这是一段完整的“计数变代数”：三个位置动作，得到一张矩阵；矩阵的平方给逆；两返回比较给倍率；再比较换块前后得到纯剪切。

\section{为什么偏要选择迹零矩阵}
二阶矩阵满足特征多项式恒等式
\[J^2-(\operatorname{tr}J)J+\det(J)I=0.
\]
因此迹为零时，$J^2=-\det(J)I$。只要行列式非零，重复同一返回就给射影逆。这比先搜索一个很长周期容易得多。

一般公式的第一项工作，就是把输出计数安排得使迹为零。第二项工作，是让两张反射只差一个简单方向。最后再证明所有取用数真的在库存内。你可以带着刚算过的 $(5,1)$ 例子，读下面每个参数。
\InsertProof{prop:odd}

\section{为什么先绕到n的三次方，再回到平方}
证明第三步有个看似绕远的地方：先得到 $\Gamma(n^3,R)$，才取得全部 $L(nR)$，最后得到 $\Gamma(n^2,R)$。

我们起初知道的只有整数步长 $L(nk)$，还不能直接把 $k$ 换成所有环元素 $t$。但每个 $t\in R$ 都能选整数 $k$ 使 $t-k\in n^2R$。所以
\[L(nt)=L(nk)L(n(t-k)).
\]
第一项已经能做，第二项属于已经取得的深子群。这样就补齐所有环参数，再应用两根定理回到平方层。每一次群的使用都有更早的真实操作作起点。

终端已经在上一讲算过：十一元五平均送到 $(-2^5,3^5,-5)$ 的方向后，三份 $-2$ 加两份 $3$ 平均成零，剩下交给含零收尾。大段矩阵最终服务的仍是这一个真实选择。
